\FloatBarrier
\section{Supplementary comparison with a single classifier}
\label{sec:singlebaseline}
\subsection{Question and comparison protocol}
A supplementary experiment tests whether a simpler classifier can reproduce the selected hybrid's balance of family detection and background errors. It was designed after the primary hybrid results were known. The comparison reuses the final benchmark's training, calibration and evaluation collections, while preserving the selected hybrid models. It tests what the compound classifier and alarm path add when the representation is retained [\evid{E9}].

Each control contains one LightGBM binary classifier trained jointly on the five attack families. The normal reference is retained, so ``single classifier'' describes the detection stage, rather than a completely unfitted preprocessing pipeline. The controls use all 116 causal features in Modena and the same features plus the 42 received-pressure-history columns in L-Town. A zero-hour control uses those causal inputs alone. A three-hour control also receives the 59 bounded-window features used by the delayed heads, giving 175 and 217 columns, respectively. Temporal features are computed from complete received series before training endpoints are sampled. Neither control receives expert predictions, routing outputs or family labels at inference.

The three-hour control is the principal complexity comparison because it shares the hybrid's maximum information allowance. It finalizes each endpoint with one delayed score and threshold; the hybrid also retains General's immediate path. The observation types and maximum deadline are aligned, while the hybrid has learned intermediate scores and aggregated evidence. Consequently, the comparison assesses the complete designs under a common allowance; it is not an isolated ablation of the router.

All ten model seeds are fitted for both controls in both networks. Each fit retains every positive TRAIN endpoint and samples up to 60,000 negative endpoints using the existing General sampling convention. The loss gives half its weight to positive endpoints, equally across families and events, and half to negatives, equally across sampled scenarios. The two controls share the sampled rows and weights within a seed. The fixed LightGBM recipe uses 400 trees, 31 leaves, learning rate 0.03, minimum leaf size 20, L1/L2 penalties of 1 and 20, and feature fraction 0.85. These settings follow the existing full temporal classifier recipe; no hyperparameter search is added.

Calibration searches the distinct score thresholds under the existing clean-FPR ceiling of 0.005. Two operating points are mandatory: maximum pooled F1, and maximum worst-family F1. Each breaks ties using the other F1 criterion and then lower clean FPR. The pooled objective defines the primary comparison; the family-balanced objective checks whether its result reflects a threshold trade-off. These single-threshold rules are explicit alternatives to the hybrid's recorded integration and protection searches.

All 40 classifiers and their 80 operating points are frozen before any supplementary evaluation. Each is then assessed on the same six source datasets used for its network's final hybrid result. Means and paired differences use the 60 model/source cells, following Equations~\ref{eq:cellmean} and~\ref{eq:paireddifference}. Both time allowances and both calibration objectives are reported. Reusing previously assessed sources makes this a supplementary controlled comparison, rather than a fresh independent confirmation.

\clearpage
\subsection{Performance and complexity trade-offs}
\begin{table}[!htbp]
\centering\small
\caption[Single-classifier comparison on Modena]{Modena: single-classifier controls versus the frozen hybrid, each averaged over ten model seeds and six shared sources. Single-classifier thresholds maximize calibration pooled F1 under the common clean-FPR ceiling; the hybrid retains its recorded operating points. Clean FPR is a percentage.}
\label{tab:singlemodena}
\begin{tabular}{@{}lrrr@{}}
\toprule
Metric & Single, 0 h & Single, 3 h & Hybrid\\
\midrule
Random F1 & 0.9672 & 0.9740 & 0.9549\\
Replay F1 & 0.9278 & 0.9060 & 0.9283\\
Drift F1 & 0.7015 & 0.7627 & 0.8274\\
Noise F1 & 0.6135 & 0.6742 & 0.8439\\
Targeted F1 & 0.9731 & 0.9795 & 0.9680\\
Family macro F1 & 0.8366 & 0.8593 & 0.9045\\
Pooled F1 & 0.8259 & 0.8509 & 0.8669\\
Clean FPR (\%) & 0.0374 & 0.0317 & 0.0584\\
\bottomrule
\end{tabular}
\end{table}
\FloatBarrier
In Modena (Table~\ref{tab:singlemodena}), the three-hour single classifier reaches pooled F1 of 0.8509, compared with 0.8669 for the hybrid. Its drift and noise means are 0.7627 and 0.6742; the hybrid reaches 0.8274 and 0.8439. Replay also improves from 0.9060 to 0.9283 in the hybrid. These gains raise family macro F1 from 0.8593 to 0.9045 and retain every hybrid family mean above 0.80. The simpler classifier has its own advantages: higher random and targeted means, and a clean FPR of 0.0317\%, compared with 0.0584\% for the hybrid. The benefit of the additional structure is stronger temporal-family coverage and pooled performance, with a higher clean-alarm rate at these operating points.

Family-balanced calibration exposes the single classifier's threshold trade-off in Modena. At three hours it raises drift and noise to 0.8311 and 0.8141, but replay falls to 0.7772 and pooled F1 to 0.7787; clean FPR rises to 0.1921\%. Thus the two prespecified single-threshold objectives recover different parts of the family profile. The hybrid preserves replay and pooled performance while retaining strong drift/noise means. Table~\ref{tab:balancedbaseline} supplies all results for the alternative calibration objective.

\begin{table}[!htbp]
\centering\small
\caption[Single-classifier comparison on L-Town]{L-Town: single-classifier controls versus the frozen hybrid, each averaged over ten model seeds and six shared sources. Single-classifier thresholds maximize calibration pooled F1 under the common clean-FPR ceiling; the hybrid retains its recorded operating points. Clean FPR is a percentage.}
\label{tab:singleltown}
\begin{tabular}{@{}lrrr@{}}
\toprule
Metric & Single, 0 h & Single, 3 h & Hybrid\\
\midrule
Random F1 & 0.9927 & 0.9950 & 0.9920\\
Replay F1 & 0.6927 & 0.7084 & 0.7312\\
Drift F1 & 0.7949 & 0.8698 & 0.8990\\
Noise F1 & 0.7706 & 0.8422 & 0.9279\\
Targeted F1 & 0.9935 & 0.9961 & 0.9918\\
Family macro F1 & 0.8489 & 0.8823 & 0.9083\\
Pooled F1 & 0.8685 & 0.8977 & 0.8950\\
Clean FPR (\%) & 0.0192 & 0.0184 & 0.0550\\
\bottomrule
\end{tabular}
\end{table}
\FloatBarrier
L-Town shows a different balance (Table~\ref{tab:singleltown}). The three-hour single classifier reaches pooled F1 of 0.8977, slightly above the hybrid's 0.8950, and a lower clean FPR of 0.0184\% versus 0.0550\%. The hybrid is stronger on replay, drift and noise: 0.7312 versus 0.7084, 0.8990 versus 0.8698, and 0.9279 versus 0.8422, respectively. These differences increase family macro F1 from 0.8823 to 0.9083. Random and targeted favor the single classifier by 0.0030 and 0.0043. The hybrid detects the weaker families better but produces more background alarms; pooled F1 is similar.

The family-balanced L-Town threshold raises the single classifier's drift and noise means to 0.8800 and 0.8653. Replay remains at 0.7085, pooled F1 is 0.8929, and clean FPR is 0.0325\%. It narrows some differences but retains lower replay, drift and noise scores than the hybrid. Across the two frozen threshold objectives, the single classifier remains an option when pooled F1 and fewer false alarms matter most. The hybrid has better coverage of replay, drift and noise.

\figfile{single_model_paired.pdf}{1.0}{Paired hybrid effects relative to the three-hour single classifier}{Family-F1 differences between the frozen hybrid and the three-hour single classifier at its pooled-F1 operating point. Each circle averages ten model seeds within one evaluation source; diamonds mark the six-source mean. Positive values favor the hybrid. The source spread is descriptive, not a confidence interval.}{fig:singlepaired}
\FloatBarrier
The paired results support these distinctions across fitted realizations. Relative to the primary three-hour control, the hybrid has higher drift, noise and family macro F1 in every one of the 60 cells in each network. Replay is higher in 54 Modena cells and all 60 L-Town cells. Pooled F1 is higher in 59 Modena cells but only 20 L-Town cells. Clean FPR is higher for the hybrid in every cell of both networks. Figure~\ref{fig:singlepaired} shows the family differences after averaging within sources; the accompanying frozen summary also records model and source variation.

Giving the single classifier three hours also improves its own pooled mean, from 0.8259 to 0.8509 in Modena and from 0.8685 to 0.8977 in L-Town. Drift and noise improve in both networks. Later evidence therefore helps a simple joint classifier as well as specialized branches. Under that allowance, the compound design improves temporal-family coverage, especially in Modena, although it does not lead on every metric. The zero- and three-hour controls are separately fitted and calibrated; their difference concerns the complete information-allowance configurations.


